% =============================================================
%  Long Technical Handbook --- Medical Image Computing (SoC)
%
%  This is the "big-picture first" handbook. For each week it gives
%  a moderately technical overview of every term the mentee will
%  encounter, with hyperlinks to deeper reads. Mentees should read
%  the relevant chapter end-to-end BEFORE opening any of the linked
%  resources or the slide decks.
%
%  Compile with: pdflatex handbook_main.tex (twice; once for ToC)
% =============================================================
\documentclass[11pt,a4paper]{article}

\usepackage[margin=2.2cm]{geometry}
\usepackage[utf8]{inputenc}
\usepackage[T1]{fontenc}
\usepackage{lmodern}
\usepackage{microtype}
\usepackage{amsmath,amssymb,amsthm,amsfonts}
\usepackage{xcolor}
\usepackage{enumitem}
\usepackage{booktabs}
\usepackage{titlesec}
\usepackage{parskip}
\usepackage[hidelinks,colorlinks=true,linkcolor=blue!50!black,urlcolor=blue!55!black,citecolor=blue!50!black]{hyperref}
\usepackage{tcolorbox}
\tcbuselibrary{breakable, skins}

\titleformat{\section}{\Large\bfseries\sffamily\color{blue!45!black}}{\thesection}{0.6em}{}
\titleformat{\subsection}{\large\bfseries\sffamily\color{blue!35!black}}{\thesubsection}{0.5em}{}
\titleformat{\subsubsection}{\normalsize\bfseries\sffamily}{\thesubsubsection}{0.4em}{}

\newtcolorbox{howtoread}{
  colback=yellow!8, colframe=orange!60!black, boxrule=0.6pt,
  arc=2pt, left=8pt,right=8pt,top=6pt,bottom=6pt, breakable,
  title=\textsf{\bfseries How to read this chapter}, fonttitle=\sffamily\bfseries
}

\newtcolorbox{keyterm}[1]{
  colback=blue!4, colframe=blue!40!black, boxrule=0.5pt,
  arc=2pt, left=8pt,right=8pt,top=4pt,bottom=4pt, breakable,
  title=\textsf{\bfseries Key idea: #1}, fonttitle=\sffamily\bfseries
}

\newcommand{\E}{\mathbb{E}}
\newcommand{\R}{\mathbb{R}}
\newcommand{\KL}{\mathrm{KL}}

\title{\sffamily\bfseries Medical Image Computing\\[3pt]
       \large A Long Technical Handbook --- Summer of Code mentee guide}
\author{\sffamily Mentor's handbook (v1)}
\date{}

\begin{document}
\maketitle

\begin{howtoread}
This handbook is your \textbf{big-picture-first} reference. Each
weekly chapter sketches every term you will encounter that week
\emph{in moderately technical language}, with hyperlinks to longer
reads. The pattern of reading you should adopt every week is:

\begin{enumerate}[nosep, leftmargin=*]
  \item \textbf{Read this chapter end-to-end.} Roughly 15--30 minutes.
        Goal: when you next hear ``Cauchy--Riemann'' or ``Hammersley--Clifford''
        or ``ELBO'', you can locate it on your mental map.
  \item \textbf{Open the week's slide deck} (see
        \texttt{slides/SLIDES\_INDEX.md}). Read it linearly; you will
        recognise the terms.
  \item \textbf{Pick at most two} linked resources per week and read
        them properly. The rest are there if you decide to specialise
        later.
  \item \textbf{Do the notebooks and exercises} in
        \texttt{week\_N/Readme.md}.
\end{enumerate}

Why this order? Because most students lose more time hopping between
random search results trying to figure out what a term means than they
do actually learning the concept. The handbook exists to give you a
\emph{scaffold} so that every external resource you read already has a
slot to drop into.
\end{howtoread}

\bigskip
\noindent\textsf{\bfseries How the repo is organised (folder $\leftrightarrow$ chapter).}
The week \emph{folders} are named by their calendar span, and each now ships its
own focused theory PDF in \texttt{week\_*/theory/}; this handbook remains the
big-picture overview across all of them.

\begin{center}
\begin{tabular}{@{}lll@{}}
\hline
Calendar & Folder & Topic \\
\hline
Weeks 1--1.5 & \texttt{week\_0\_to\_1.5} & Images, math \& probability foundations \\
Weeks 1.5--3 & \texttt{week\_1.5\_to\_3} & Statistical priors, denoising, inverse problems \\
Week 3       & \texttt{week\_3}          & CT \& MRI physics (fundamentals) \\
Week 4       & \texttt{week\_4}          & Shape, segmentation, registration, kernels \\
Weeks 5--6.5 & \texttt{week\_5}, \texttt{week\_6} & CNNs; VAEs and GANs \\
Weeks 6.5--8 & \texttt{week\_7}, \texttt{week\_8} & Paper reading; final project \\
\hline
\end{tabular}
\end{center}

\noindent Where a chapter below says ``Week 2'' or ``Week 3'', read it as the
\emph{calendar} label in this table. The reconstruction / inverse-problems
material lives in \texttt{week\_1.5\_to\_3}; the CT and MRI \emph{physics} moved
to their own \texttt{week\_3} chapter.

\tableofcontents
\newpage

% =============================================================
\section{Course philosophy: math first, then DL}
% =============================================================

Medical Image Computing (MIC) is the engineering and mathematical
discipline that builds algorithms for medical-image data. It sits at
the intersection of computer science, applied mathematics, signal /
image processing, statistics, and the physics of imaging modalities
(X-rays, magnetic resonance, ultrasound, microscopy, $\ldots$).

This course is built on top of \emph{CS-736 Medical Image Computing} at
IIT Bombay (Prof.~Suyash P.~Awate). Two emphases follow from that:

\begin{itemize}[leftmargin=*,nosep]
  \item It is \textbf{math-first}. Every algorithm in weeks 2--6 is
        either \emph{a maximum-a-posteriori (MAP) estimate under some
        prior and noise model}, or \emph{a clustering / classification
        / regression problem with a chosen geometry}. Deep learning is
        a powerful \emph{parameterisation}, but it does not replace
        the principled scaffold.
  \item It is \textbf{modality-aware}. CT and MRI images are produced
        by physical systems with very specific structure; the
        algorithms that work on them inherit that structure
        (Beer--Lambert law, Bloch equations, Radon transform, $k$-space).
\end{itemize}

The recurring object in everything we do is Bayes' theorem applied to
an unknown image (or label / shape / mapping) $x$ given measurements $y$:
\begin{equation}
  p(x \mid y) \;\propto\; \underbrace{p(y \mid x)}_{\text{noise / forward model}} \cdot \underbrace{p(x)}_{\text{prior}}.
  \label{eq:bayes-image}
\end{equation}

Almost everything you learn this summer is one of:

\begin{itemize}[leftmargin=*,nosep]
  \item a way to write down $p(y \mid x)$ (Weeks 2, 3);
  \item a way to write down $p(x)$ (Weeks 2, 4, 5, 6);
  \item a way to maximise the right-hand side of \eqref{eq:bayes-image}
        (Weeks 3, 5);
  \item a way to evaluate whether your $\hat x$ is any good (Weeks 5, 7, 8).
\end{itemize}

% =============================================================
\section{Weeks 1--1.5 (\texttt{week\_0\_to\_1.5}) --- Images, math \& probability foundations}
% =============================================================

\subsection{Big picture}

A \textbf{digital image} is just an $N$-dimensional array of numbers
laid out on a regular grid. A pixel (2D) or voxel (3D) is one entry
in that array; its \emph{intensity} is the number stored there. What
that number physically means depends on the modality: for CT it is an
attenuation coefficient (or its Hounsfield-unit rescaling); for MRI it
is the magnitude of a complex magnetisation; for microscopy it is a
photon count.

The colour you \emph{see} when an image is displayed is not the data ---
it is a \textbf{colormap}, a function from the data values to colours.
The same MRI slice can look black-on-white or white-on-black; the data
are identical. Get into the habit of always inspecting the
\emph{intensity histogram} of an image alongside the picture.

Two pieces of mathematical machinery sit underneath nearly every
algorithm in this course:

\begin{itemize}[leftmargin=*,nosep]
  \item \textbf{Linear algebra} --- because images are vectors (just
        very long ones), and almost every operation we do is a linear
        map followed by something. Make sure you are comfortable with:
        matrix--vector multiplication as a linear map, transposes,
        inverses, ranks, eigen-decomposition of symmetric matrices,
        the singular value decomposition (SVD) at the level of
        ``it factorises $A = U\Sigma V^\top$ and the columns of $V$
        are the input principal directions''.
  \item \textbf{Probability \& statistics} --- because every imaging
        system produces noisy measurements, and every prior we
        impose on $x$ is a probability distribution. You should be
        able to write the PDF of a Gaussian and a Poisson by hand,
        and state Bayes' theorem without looking it up.
\end{itemize}

\subsection{Vocabulary you will meet}

\begin{description}[leftmargin=*]
  \item[\textbf{Pixel / voxel.}] A single sampling location of the
        image grid. ``Pix-el'' = \emph{pic}ture \emph{el}ement. Voxel
        is the 3D analogue.
  \item[\textbf{Intensity.}] The value stored at a pixel. Real or
        integer; signed or unsigned; range depends on bit-depth
        ($[0,255]$ for 8-bit, $[0,65535]$ for 16-bit, real for float
        scanners).
  \item[\textbf{Colormap.}] A function from intensity to displayed
        colour. Examples: \texttt{gray}, \texttt{viridis}, hot-iron.
  \item[\textbf{Voxel spacing.}] The physical size (in mm) of one
        voxel along each axis. Medical scanners produce
        \emph{anisotropic} voxels (e.g.\ $1{\times}1{\times}3$ mm)
        for many protocols.
  \item[\textbf{Convolution.}] The operation
        $(I * k)[m,n] = \sum_{i,j} I[m-i,n-j]\,k[i,j]$. Equivalent in
        spectrum: pointwise product in the Fourier domain.
  \item[\textbf{2D Fourier transform.}] Decomposes an image into a
        sum of plane waves indexed by spatial frequency $(u,v)$.
  \item[\textbf{Norms.}] $\ell^2$-norm
        $\|x\|_2 = \sqrt{\sum_i x_i^2}$, $\ell^1$
        $\|x\|_1 = \sum_i |x_i|$, $\ell^\infty$
        $\|x\|_\infty = \max_i |x_i|$. The $\ell^1$ norm promotes
        sparsity --- a fact we will use in Week 3.
  \item[\textbf{Inner product.}] $\langle x,y\rangle = x^\top y$ in
        Euclidean space. Norms are induced from inner products via
        $\|x\| = \sqrt{\langle x,x\rangle}$.
  \item[\textbf{Eigenvector / eigenvalue.}] $Av = \lambda v$. For
        symmetric $A$, eigenvectors are orthogonal and eigenvalues
        are real; we will exploit this repeatedly (PCA, kPCA, MRF
        analysis).
  \item[\textbf{Positive (semi)definite matrix.}] $x^\top A x \ge 0$
        for all $x$. Covariance matrices are PSD. Required as kernel
        matrices in Week 4.
  \item[\textbf{Random variable, PDF, CDF.}] If $X$ has PDF $p_X$,
        then $\Pr(X \le a) = \int_{-\infty}^a p_X(t)\,dt$.
  \item[\textbf{Gaussian distribution.}]
        $\mathcal{N}(\mu,\sigma^2)$: $p(x) = \tfrac{1}{\sqrt{2\pi}\sigma}
         \exp\!\big(-\tfrac{(x-\mu)^2}{2\sigma^2}\big)$.
  \item[\textbf{Poisson distribution.}]
        $\Pr(X = k) = \tfrac{\lambda^k e^{-\lambda}}{k!}$.
        Mean = variance = $\lambda$. Models photon counting.
  \item[\textbf{Bayes' theorem.}]
        $p(x \mid y) = p(y \mid x)\,p(x) / p(y)$. The denominator
        $p(y)$ does not depend on $x$ and is usually ignored when we
        \emph{maximise} the posterior over $x$.
  \item[\textbf{MLE vs.\ MAP.}] Maximum-likelihood estimate:
        $\hat x_{\mathrm{MLE}} = \arg\max_x p(y \mid x)$. Maximum
        a-posteriori: $\hat x_{\mathrm{MAP}} = \arg\max_x p(y \mid x)\,p(x)$.
        MAP $=$ MLE plus a prior.
  \item[\textbf{Gradient descent.}] $x_{k+1} = x_k - \eta \nabla f(x_k)$,
        $\eta > 0$ a step size.
\end{description}

\subsection{Pointers}

3Blue1Brown's \emph{Essence of Linear Algebra} is the right intuition
playlist; StatQuest's videos on probability vs.\ likelihood and on
Bayes' theorem are the right intuition for probability. For depth, MIT
OCW 18.06 (Strang) and 6.041 (Tsitsiklis) are still the gold standard.
The full per-week resource list is in \texttt{week\_1/Readme.md}.

% =============================================================
\section{Weeks 1.5--3 (\texttt{week\_1.5\_to\_3}), part 1 --- Statistical priors and image denoising}
% =============================================================

\subsection{Big picture}

Week 2 is the first time we use \eqref{eq:bayes-image} for real. The
problem is \emph{denoising}: given a noisy image $y$, recover a
clean image $\hat x$.

We need two ingredients:
\begin{enumerate}[nosep, leftmargin=*]
  \item A \textbf{noise model} $p(y \mid x)$, that depends on the
        modality. Gaussian for thermal noise; Poisson for
        photon-counting (X-ray, CT, microscopy); Rician for the
        magnitude of an MRI image; Compound-Poisson for
        polychromatic CT; Rician/Rayleigh/K for ultrasound.
  \item A \textbf{prior} $p(x)$ that encodes what a ``plausible''
        image looks like --- usually a mixture of (i) intensities vary
        smoothly inside organs and (ii) large jumps are allowed only
        across organ boundaries.
\end{enumerate}

The right mathematical object for the prior is a \textbf{Markov
Random Field (MRF)}. An MRF is a probability distribution over an
image (treated as a vector of $N$ random variables, one per pixel)
in which each pixel is conditionally independent of all distant
pixels given its neighbours. By the \textbf{Hammersley--Clifford
theorem}, every (positive) MRF can be written as a
\textbf{Gibbs distribution},
\begin{equation}
  p(x) = \tfrac{1}{Z}\exp\!\Big(-\tfrac{1}{T}\sum_{c \in C} V_c(x_c)\Big),
  \label{eq:gibbs}
\end{equation}
where the sum is over \emph{cliques} $c$ (subsets of pixels that are
all pairwise neighbours), $V_c$ are \emph{clique potentials} you get
to design, $T$ is a temperature, and $Z$ is the normalising
\emph{partition function}.

Different choices of clique potential give different priors. A
quadratic potential $V(x_i, x_j) = (x_i - x_j)^2$ gives a
\textbf{Gaussian MRF}; it strongly penalises any intensity jump and
therefore blurs edges. \textbf{Edge-preserving} potentials
(Huber, generalised Gaussian, truncated quadratic, Mumford--Shah) cap
the penalty for large jumps and so preserve organ boundaries.
\textbf{Ising} (binary labels) and \textbf{Potts} (multi-label) priors
do the analogous job for segmentation, which we will use in Week 4.

Putting it together, \textbf{Bayesian denoising = MAP estimation}:
\begin{equation}
  \hat x = \arg\max_x p(y \mid x)\,p(x)
         = \arg\min_x \big[\underbrace{-\log p(y \mid x)}_{\text{data fidelity}} \;+\;
                     \underbrace{-\log p(x)}_{\text{regularisation}}\big].
\end{equation}

The negative-log Gibbs prior \emph{is} the regulariser you have
already seen in machine learning. A Gaussian prior gives $\ell^2$
(Tikhonov / ridge) regularisation; a Laplace prior on pixel
differences gives \emph{total variation (TV)} ($\ell^1$); an MRF with
edge-preserving potentials gives a robust regulariser.

\subsection{Vocabulary you will meet}

\begin{description}[leftmargin=*]
  \item[\textbf{Random field.}] A collection of random variables
        $\{X_i\}$ indexed by sites $i$ on a grid.
  \item[\textbf{Neighbourhood system.}] For each site $i$, a set
        $N_i$ of neighbouring sites. 4-/8-neighbour in 2D;
        6-/18-/26-neighbour in 3D.
  \item[\textbf{Clique.}] A subset of sites that are all pairwise
        neighbours (so for a 4-neighbour 2D system, cliques are
        singletons and edges only).
  \item[\textbf{Markov property.}] $p(x_i \mid x_{S\setminus\{i\}}) =
        p(x_i \mid x_{N_i})$.
  \item[\textbf{Hammersley--Clifford theorem (1971).}] Under
        positivity ($p(x) > 0$), an MRF and a Gibbs distribution on
        the same neighbourhood system are the same object.
  \item[\textbf{Clique potential $V_c$.}] A scalar function over the
        values of pixels in a clique $c$; you design these to encode
        ``smoothness'' or ``sparsity'' or ``label agreement''.
  \item[\textbf{Partition function $Z$.}] The normaliser in
        \eqref{eq:gibbs}. Almost always intractable to compute in
        closed form; we can still optimise the energy without ever
        evaluating $Z$.
  \item[\textbf{Ising / Potts model.}] Binary / multi-label MRFs
        whose clique potential is $V(x_i,x_j) = \mathbb{1}[x_i \neq x_j]$.
  \item[\textbf{Gaussian MRF.}] Quadratic potentials; closed-form
        posterior; smooths edges.
  \item[\textbf{Edge-preserving prior.}] Potential that grows slowly
        (or saturates) for large differences. Huber:
        $V(t) = t^2/2$ for $|t|\le\delta$ and $\delta|t| - \delta^2/2$
        otherwise.
  \item[\textbf{Gaussian noise.}] $y_i = x_i + \epsilon_i$ with
        $\epsilon_i \sim \mathcal{N}(0,\sigma^2)$, i.i.d.
  \item[\textbf{Poisson noise.}] $y_i \sim \mathrm{Poisson}(\lambda_i)$
        with $\lambda_i$ a function of the underlying intensity $x_i$
        and the imaging geometry.
  \item[\textbf{Compound Poisson.}] Sum of $N$ i.i.d.\ random
        variables where $N$ itself is Poisson; arises in
        polychromatic CT.
  \item[\textbf{Rician distribution.}] Distribution of the magnitude
        $\sqrt{a^2 + b^2}$ when $(a,b) \sim \mathcal{N}((\nu\cos\phi,\nu\sin\phi),\sigma^2 I)$.
        Its PDF involves the modified Bessel function $I_0$. Used for
        MRI magnitude images. Reduces to Rayleigh when $\nu = 0$.
  \item[\textbf{Signal-to-noise ratio.}] $\mathrm{SNR} = \E[X] / \mathrm{SD}(X)$.
        For Poisson, $\mathrm{SNR} = \sqrt{\lambda}$.
  \item[\textbf{Total variation (TV).}] Regulariser
        $\mathrm{TV}(x) = \sum_i \|\nabla x\|_2$ (isotropic) or
        $\sum_i (|x_{i+1}-x_i| + |x_{i+\text{row}}-x_i|)$ (anisotropic).
        Promotes piecewise-constant images.
  \item[\textbf{Non-local means.}] A denoising algorithm that
        averages each pixel with all pixels whose \emph{patch} is
        similar; not strictly MAP, but a beautifully simple
        algorithm rooted in self-similarity.
\end{description}

\subsection{Pointers}

Bouman's free PDF \emph{Model-Based Image Processing} is the natural
textbook for this week and Week 3. Geman \& Geman (1984) and Buades et
al.\ (2005) are the two classical papers to read if you want depth.
The Gudbjartsson \& Patz (1995) Rician note is a beautiful 4-page
read.

% =============================================================
\section{Weeks 1.5--3 \& Week 3 --- Inverse problems (\texttt{week\_1.5\_to\_3}) and CT/MRI physics (\texttt{week\_3})}
% =============================================================

\subsection{Big picture}

In Week 2 we recovered $\hat x$ from a noisy version of itself.
In Week 3 we make it harder: we observe $y$ that has gone through a
\emph{forward operator} $G$ that scrambles the image:
\begin{equation}
  y = Gx + n.
  \label{eq:fwd}
\end{equation}
This is called an \textbf{inverse problem}. In CT, $G$ is the discrete
Radon transform (sum of all line integrals through the patient); in
MRI, $G$ is a (typically under-sampled) 2D Fourier transform. In both
cases $G$ may be \emph{rank-deficient} (you don't get enough
measurements to invert) or \emph{ill-conditioned} (small noise in
$y$ blows up in $\hat x$). This is what ``\textbf{ill-posed}'' means.

\textbf{MAP reconstruction} fixes ill-posedness by adding a prior:
\begin{equation}
  \hat x = \arg\min_x \;\;\tfrac{1}{2\sigma^2}\|Gx - y\|_2^2 \;+\; \lambda \, R(x),
  \label{eq:map-recon}
\end{equation}
where $R(x)$ is the negative-log Week-2 prior (Tikhonov,
edge-preserving MRF, TV, dictionary, $\ldots$) and $\lambda$
trades data-fidelity against the prior.

To solve \eqref{eq:map-recon} you need
\textbf{multivariate gradient descent} (and friends: Newton,
quasi-Newton, ADMM, primal-dual). The gradient of the data term is
\begin{equation}
  \nabla_x \tfrac{1}{2}\|Gx - y\|_2^2 \;=\; G^\top (Gx - y).
\end{equation}
Derive it once by hand using the chain rule with
$w := Gx - y, f := \tfrac12 w^\top w$. For MRI, $x$ is
complex-valued, and you need to be careful with what
``\emph{gradient of a real-valued function of a complex argument}''
means; this is the only place the
\textbf{Cauchy--Riemann equations} show up in this course.

\subsection{CT physics}

X-rays are short-wavelength electromagnetic radiation. When an X-ray
of intensity $I_0$ passes through a thin slab of tissue of thickness
$\Delta x$ with \textbf{attenuation coefficient} $A(x)$, a fraction
$A(x)\Delta x$ of the photons get absorbed. This linear relationship,
in the infinitesimal limit, is the
\textbf{Beer--Lambert law}:
\begin{equation}
  I_1 \;=\; I_0 \,\exp\!\Big(-\!\int_{x_0}^{x_1} A(x)\,dx\Big).
  \label{eq:beer}
\end{equation}
A single X-ray therefore gives you a \emph{line integral} of $A$. One
projection is not enough to recover $A(\cdot)$; to do tomography,
the scanner takes many such projections from many angles. The
collection of all line integrals at all angles is the
\textbf{Radon transform} of the image.

The \textbf{Fourier slice theorem} says that the 1D Fourier transform
of a 1D parallel projection at angle $\theta$ equals a 1D radial
slice through the 2D Fourier transform of the image at the same
angle. This is the basis of \textbf{filtered back-projection (FBP)} ---
the direct, non-iterative CT reconstruction algorithm. Iterative
\textbf{algebraic reconstruction (ART)} and MAP estimation are the
alternatives, used when the data are sparse-view, low-dose, or
metal-corrupted.

\subsection{MRI physics}

The human body is mostly water, and water is rich in $^1$H
hydrogen nuclei (protons). A spinning proton has a magnetic dipole
moment; placed in a strong static magnetic field $B_0$, the dipoles
align and \emph{precess} around $B_0$ at the
\textbf{Larmor frequency} $\omega_0 = \gamma B_0$, where $\gamma$ is
the gyromagnetic ratio. For $^1$H, $\gamma \approx 42.58$ MHz/T, so
at 3 T you have $\omega_0 \approx 128$ MHz.

The dynamics of the aggregate magnetisation
$\mathbf{M}(t) = (m_x, m_y, m_z)$ are described by the
\textbf{Bloch equations}: a system of ODEs containing three terms ---
precession, longitudinal relaxation toward $m_{eq}$ with time
constant $T_1$, and transverse relaxation toward $0$ with time
constant $T_2$. $T_1$ and $T_2$ depend on tissue, which is why MRI
can distinguish soft tissues. A short RF pulse tips
the magnetisation into the transverse plane; gradient fields encode
spatial position into the precession frequency
(\emph{frequency encoding}) and phase
(\emph{phase encoding}). Each MRI measurement is one complex sample
in \textbf{$k$-space} (the 2D/3D Fourier domain of the image). An
\emph{inverse FFT} of a complete $k$-space dataset gives the image.

When we sub-sample $k$-space (to image faster), the inverse FFT
aliases; the cure is to combine the sparse measurement with a prior
that gives a unique solution --- the
\textbf{compressed-sensing MRI} story that you will meet at the end
of the week.

\subsection{Vocabulary you will meet}

\begin{description}[leftmargin=*]
  \item[\textbf{Forward operator $G$.}] The matrix that maps the
        unknown image to the measurement. Radon transform for CT;
        Fourier transform for MRI.
  \item[\textbf{Ill-posed.}] Either no solution exists, the solution
        is not unique, or it does not depend continuously on the
        data (Hadamard's definition).
  \item[\textbf{Tikhonov regularisation.}] $R(x) = \|x\|_2^2$ (or
        $\|Lx\|_2^2$ for some linear $L$, e.g.\ a gradient).
  \item[\textbf{Total variation (TV).}] $R(x) = \|\nabla x\|_1$. The
        $\ell^1$ on gradients makes the prior sparse on edges.
  \item[\textbf{Line search.}] Pick the step size $\eta$ in gradient
        descent by minimising $f(x_k - \eta g_k)$ along the ray.
        Variants: backtracking, golden-section, Newton's method.
  \item[\textbf{Matrix calculus.}] The Jacobian
        $\partial f / \partial x$ has $m$ rows (outputs) and $n$
        columns (inputs); the gradient of a scalar function is the
        transpose of the Jacobian (so it lives in the same space as
        $x$).
  \item[\textbf{Cauchy--Riemann equations.}] $u_x = v_y$,
        $u_y = -v_x$ for $f(x+iy) = u(x,y) + iv(x,y)$ to be
        complex-differentiable. We invoke them only to justify
        treating complex MRI gradients via the
        ``\emph{complex-as-doubled-real}'' trick.
  \item[\textbf{Attenuation coefficient.}] Local rate of photon
        absorption per unit length. Units: $1/\text{length}$.
  \item[\textbf{Hounsfield unit.}] Affine rescaling of attenuation
        coefficient with water at $0$ HU and air at $-1000$ HU.
  \item[\textbf{Beer--Lambert law.}] Eq.~\eqref{eq:beer}.
  \item[\textbf{Radon transform.}] Collection of all parallel
        line-integrals through an image at all angles.
  \item[\textbf{Fourier slice theorem.}] Links 1D projections to 2D
        Fourier slices; the theoretical basis of FBP.
  \item[\textbf{Filtered back-projection (FBP).}] Closed-form CT
        reconstruction by applying a ramp filter to each projection,
        then back-projecting.
  \item[\textbf{ART (algebraic reconstruction).}] Iteratively project
        the current estimate onto each measurement constraint;
        converges to a least-norm solution.
  \item[\textbf{Larmor frequency.}] $\omega_0 = \gamma B_0$.
  \item[\textbf{Bloch equations.}] ODEs for the magnetisation vector;
        contain $T_1$ and $T_2$ relaxation times.
  \item[\textbf{$T_1$ / $T_2$.}] Tissue-dependent relaxation times.
        $T_1$ governs longitudinal recovery, $T_2$ transverse decay.
  \item[\textbf{$k$-space.}] The 2D/3D Fourier domain of the image;
        what the MRI scanner actually samples.
  \item[\textbf{Compressed sensing.}] Recover a signal from
        sub-Nyquist samples by exploiting sparsity in some
        transform domain.
\end{description}

\subsection{Pointers}

Kak \& Slaney's \emph{Principles of Computerized Tomographic Imaging}
and Bouman's \emph{MBIP} are the two free textbooks to keep at hand.
For MRI, watch any one of Allen Ye's videos and read the iMaios
$k$-space explainer; for compressed sensing read the Lustig--Donoho
SparseMRI paper.

% =============================================================
\section{Week 4 (\texttt{week\_4}) --- Shape, classical segmentation, registration, kernels}
% =============================================================

\subsection{Big picture}

We now switch from \emph{recovering an image} to \emph{recovering
structure}: which voxels belong to which tissue
(\textbf{segmentation}), how a population of anatomical structures
varies in geometry (\textbf{shape analysis}), how to align two scans
to the same coordinate frame (\textbf{registration}), and an extra
toolkit (\textbf{kernel methods}) that everything from SVMs to kernel
PCA to many shape models depends on.

\subsubsection{Segmentation}

The simplest segmenter is \textbf{$k$-means clustering}: partition
$N$ feature vectors $\{x_n\}$ into $K$ clusters by minimising
\begin{equation}
  \sum_{k=1}^{K}\sum_{n \in S_k} \|x_n - \mu_k\|_2^2.
\end{equation}
Solved by alternating optimisation: fix means, update assignments
(each point goes to nearest mean); fix assignments, update means (set
each $\mu_k$ to the centroid of its cluster). It always converges (to
a local minimum); finding the global minimum is NP-hard.
\textbf{$k$-means++} is a clever initialisation that gives a
provable approximation factor. \textbf{$k$-medoids} replaces means
by data points and gives outlier robustness. The Voronoi cells of
the means partition feature space; if you switch the distance from
Euclidean to Manhattan, you switch from straight-edge to
axis-aligned cells.

\textbf{Gaussian Mixture Models (GMMs)} are the soft-assignment
generalisation: each cluster is a Gaussian, points have a
\emph{posterior probability} of belonging to each cluster, and the
parameters are estimated by the \textbf{Expectation--Maximisation
(EM)} algorithm. EM is a coordinate-ascent algorithm on a lower
bound of the log-likelihood (the same machinery you will reuse in
Week 6 for the VAE's ELBO).

\textbf{Graph cuts} formulate segmentation as min-cut / max-flow on a
graph whose edges encode unary (label likelihood) and pairwise
(smoothness) potentials --- in other words, MAP inference on an
MRF with binary or multi-label cliques. \textbf{Spectral clustering}
uses the eigenvectors of the graph Laplacian instead. \textbf{Level
sets} (Chan--Vese, Mumford--Shah) evolve an implicit curve by a PDE
whose driving force is an energy.

\subsubsection{Shape analysis}

A \textbf{shape} is what remains of an outline after you remove
translation, rotation, and isotropic scale --- formally, an
equivalence class under the
\textbf{similarity group}. We represent a shape by an ordered
\textbf{pointset} of landmarks placed consistently across all
subjects; that fixes correspondences and lets us write each shape
as a vector in $\R^{DN}$.

To remove the nuisance transformations we use
\textbf{Procrustes analysis}: translate to the origin, scale to
unit norm, and rotate to minimise the sum-of-squared distances to a
reference shape (closed form via SVD of the cross-covariance).
\textbf{Generalised Procrustes Analysis (GPA)} aligns all shapes
to a moving mean iteratively. After alignment, the shapes live near
a single point on a high-dimensional sphere, and a
\textbf{tangent-space approximation} lets us run linear PCA on
them. The principal components are the \textbf{modes of shape
variation}; walking $\pm 2\sigma$ along each mode is the
classic visualisation. The model that does all of this is the
\textbf{Active Shape Model (ASM)}; if you add per-vertex intensity
information you get the \textbf{Active Appearance Model (AAM)}.

\subsubsection{Registration}

Given two images $R(x)$ and $T(x)$, find a transformation
$\phi: \Omega \to \Omega$ that ``aligns'' $T \circ \phi$ to $R$.
Transformations are usually rigid, similarity, affine (linear),
B-spline (parametric non-linear), or \textbf{diffeomorphic} (smooth
and invertible, the gold standard for medical work). Similarity
measures depend on modality:
\textbf{sum-of-squared differences (SSD)} for same modality,
\textbf{normalised cross-correlation (NCC)} for robustness to
linear intensity changes, \textbf{mutual information (MI)} for
\emph{different} modalities (CT $\leftrightarrow$ MRI). Atlas
building (e.g.\ Talairach for the brain) iteratively registers all
subjects to a moving mean (same recipe as GPA, but in image space).

\subsubsection{Kernels (supplementary)}

Many of the methods above (and a few in Week 5) are special cases
of a unified picture. Map data nonlinearly through some
$\phi : \R^d \to \mathcal{F}$ into a (possibly
infinite-dimensional) Hilbert space $\mathcal{F}$ and do linear
methods there. The \textbf{kernel trick} is that you never need to
construct $\phi$ explicitly --- you just need the inner product
$k(x,y) = \langle \phi(x), \phi(y) \rangle_{\mathcal{F}}$. The
classic example is the
\textbf{RBF (Gaussian) kernel} $k(x,y) = \exp(-\gamma\|x-y\|^2)$,
whose feature space is infinite-dimensional.

A kernel is admissible iff it is a positive-definite function
(\textbf{Mercer / Moore--Aronszajn}); the corresponding feature space
is a \textbf{Reproducing Kernel Hilbert Space (RKHS)}. \textbf{Kernel
PCA} runs PCA in the RKHS using only the \textbf{Gram matrix}
$K_{ij} = k(x_i, x_j)$. \textbf{Support Vector Machines (SVMs)} find
a maximum-margin hyperplane in the RKHS; deriving them goes through
the \textbf{KKT conditions} of constrained convex optimisation.

\subsection{Vocabulary you will meet}

\begin{description}[leftmargin=*]
  \item[\textbf{Cluster compactness.}] Within a cluster,
        $\sum_n \|x_n - c\|^2 = \tfrac{1}{N}\sum_{n,j}\|x_n - x_j\|^2$;
        the identity that makes $k$-means about pairwise distances
        as much as about centroids.
  \item[\textbf{Voronoi diagram.}] Partition of space by ``nearest
        centre''; the cell shapes depend on the distance.
  \item[\textbf{Expectation--Maximisation (EM).}] Alternating
        maximisation of $\E_{q(z)}[\log p(x,z;\theta)] - \KL(q\|p(z|x;\theta))$
        over $q$ then $\theta$. Same idea reappears in the
        VAE in Week 6.
  \item[\textbf{Graph Laplacian.}] $L = D - W$, where $D$ is the
        degree matrix and $W$ the affinity matrix. Eigenvectors
        give spectral embeddings.
  \item[\textbf{Min-cut / max-flow.}] Polynomial-time algorithms for
        graph cuts in undirected graphs; the workhorse of binary
        graph-cut segmentation.
  \item[\textbf{Level set.}] An evolving curve represented as the
        zero level set of an embedding function $\phi$, governed by
        a PDE $\partial \phi/\partial t = F\|\nabla\phi\|$.
  \item[\textbf{Procrustes alignment.}] Closed-form removal of
        translation, scale, rotation between two pointsets via SVD.
  \item[\textbf{Tangent space.}] Local linear approximation of a
        curved manifold (here, shape space).
  \item[\textbf{Active Shape Model.}] Mean shape $\bar x$ plus
        principal directions $\{v_k\}$:
        $x \approx \bar x + \sum_k b_k v_k$.
  \item[\textbf{Mutual information.}]
        $\mathrm{MI}(X;Y) = \KL(p_{XY}\|p_X p_Y)$. The right
        similarity measure for multi-modal registration.
  \item[\textbf{Diffeomorphism.}] Smooth invertible map with smooth
        inverse; the gold standard for non-linear medical-image
        registration.
  \item[\textbf{Kernel trick.}] Inner products in feature space via
        $k(x,y)$, never building $\phi$.
  \item[\textbf{Positive-definite function.}] $\sum_{ij} c_i c_j k(x_i, x_j) \ge 0$
        for all choices of $\{x_i\}$, $\{c_i\}$.
  \item[\textbf{RKHS / Moore--Aronszajn.}] Every PD function $k$ is
        the reproducing kernel of a unique Hilbert space of
        functions; that's where $\phi(x) := k(x,\cdot)$ lives.
  \item[\textbf{Gram matrix.}] $K_{ij} = k(x_i, x_j)$; what kernel
        PCA actually eigendecomposes.
  \item[\textbf{KKT conditions.}] First-order optimality conditions
        for constrained convex programs; derived in Boyd \& Vandenberghe
        ch.\ 5. You meet them when you derive the SVM dual.
\end{description}

\subsection{Pointers}

Bishop's \emph{Pattern Recognition and Machine Learning} chapters 9
(clustering, EM) and 12 (PCA) cover most of this. The von Luxburg
spectral-clustering tutorial is the canonical read. Cootes et al.\
(1995) is the original ASM paper. For kernels and SVMs, Sch\"olkopf
\& Smola's \emph{Learning with Kernels} and Boyd \& Vandenberghe's
\emph{Convex Optimization} (KKT in ch.\ 5) are the references.

% =============================================================
\section{Weeks 5--6.5 (\texttt{week\_5}), part 1 --- CNNs for medical images}
% =============================================================

\subsection{Big picture}

A \textbf{convolutional neural network (CNN)} is a stack of layers
each of which slides one or more small filters (``kernels'') across
its input, applies a non-linearity, possibly downsamples, and feeds
the result to the next layer. Two properties make CNNs the right
inductive bias for images:
\textbf{translation equivariance} (same filter, applied
everywhere) and \textbf{locality} (each output pixel depends only
on a neighbourhood of the input).

For \textbf{classification}, the typical recipe is
Conv$\to$Norm$\to$ReLU $\to$ Pool repeated $L$ times, then a
fully-connected head producing class scores. For
\textbf{segmentation} the classical choice in medical imaging is
the \textbf{U-Net} (Ronneberger et al., MICCAI 2015): an
encoder-decoder with skip connections that concatenate
encoder feature maps to the matching-resolution decoder layer. The
skips let the decoder recover fine detail that pooling has
destroyed.

The \textbf{engineering recipe} that wins most medical segmentation
challenges is \textbf{nnU-Net} (Isensee et al., 2021) --- a U-Net
plus a careful, dataset-adaptive pipeline (preprocessing,
augmentation, training schedule, postprocessing). The take-away is
that in medical imaging, \emph{good engineering of a standard
architecture often beats a novel architecture}.

Training a CNN means solving a non-convex optimisation problem with
\textbf{stochastic gradient descent} (and friends: Adam, AdamW).
Gradients are computed by \textbf{reverse-mode automatic
differentiation} (the ``backpropagation'' algorithm). The whole
thing is one giant chain rule on the computational graph.

\subsection{Vocabulary you will meet}

\begin{description}[leftmargin=*]
  \item[\textbf{Convolution layer.}] Slides a small kernel across the
        input; each kernel is shared across spatial positions
        (translation equivariance). Has $C_\text{in} \cdot C_\text{out}
        \cdot k \cdot k$ parameters for a $k \times k$ 2D kernel.
  \item[\textbf{Cross-correlation vs.\ convolution.}] DL libraries
        compute cross-correlation; for learning that's identical to
        true convolution because the kernel is also learnt.
  \item[\textbf{Stride / padding / dilation.}] Three knobs that
        determine output size and receptive field.
  \item[\textbf{Receptive field.}] Spatial extent of the input that
        a given output pixel depends on. Grows with depth and
        dilation.
  \item[\textbf{Activation.}] Pointwise non-linearity: ReLU
        $\max(0,x)$; GELU $x \cdot \Phi(x)$ (Gaussian CDF);
        sigmoid $1/(1+e^{-x})$; softmax (vector $\to$ probability
        simplex); softplus $\log(1+e^x)$.
  \item[\textbf{Pooling.}] Downsamples by max or average; mostly
        replaced by strided convolution in modern nets.
  \item[\textbf{Normalisation.}] BatchNorm normalises per-channel
        across a batch; GroupNorm (preferred for small medical
        batches) per group of channels; InstanceNorm per sample.
  \item[\textbf{Dice loss.}]
        $1 - \tfrac{2|P \cap T|}{|P| + |T|}$. Handles class
        imbalance better than cross-entropy for medical
        segmentation.
  \item[\textbf{Backpropagation.}] Apply the chain rule layer-by-layer
        in reverse to compute all parameter gradients in one pass.
  \item[\textbf{Adam / AdamW.}] Adaptive-learning-rate optimisers;
        \emph{AdamW} (decoupled weight decay) is the default for
        modern training.
  \item[\textbf{U-Net.}] Encoder-decoder CNN with skip connections,
        the medical-segmentation workhorse.
  \item[\textbf{Dice / IoU / Hausdorff 95 / ASSD.}] Segmentation
        metrics. \emph{Dice} is overlap; \emph{IoU} is intersection
        over union; \emph{HD95} is the 95th-percentile Hausdorff
        distance; \emph{ASSD} is average symmetric surface
        distance.
\end{description}

\subsection{Pointers}

Stanford CS231n class notes are the right text. 3Blue1Brown's first
two NN videos for intuition. The original U-Net and nnU-Net papers
are short and worth reading in full. The MONAI tutorials are the
fastest path to a working medical-imaging PyTorch pipeline.

% =============================================================
\section{Weeks 5--6.5 (\texttt{week\_6}), part 2 --- VAEs and GANs}
% =============================================================

\subsection{Big picture}

A \textbf{generative model} is a probability distribution
$p_\theta(x)$ over the data $x$ that you can both \emph{score} and
(usually) \emph{sample} from. Once you have a generative model, you
have a \emph{learned prior} that you can plug back into the
Bayesian reconstruction objective \eqref{eq:map-recon} from Week 3.

Two classical deep generative families:

\subsubsection{Variational Autoencoder (VAE)}

Assume $x$ is generated from a latent code $z$ through a decoder
$p_\theta(x \mid z)$, with a simple prior $p(z) = \mathcal{N}(0,I)$.
The marginal $p_\theta(x) = \int p_\theta(x \mid z) p(z) dz$ is
intractable. Introduce an \textbf{encoder}
$q_\phi(z \mid x) \approx p_\theta(z \mid x)$, and derive the
\textbf{Evidence Lower Bound (ELBO)}:
\begin{equation}
  \log p_\theta(x) \;\ge\; \E_{q_\phi(z|x)}[\log p_\theta(x|z)] \;-\; \KL\big(q_\phi(z|x) \,\|\, p(z)\big).
  \label{eq:elbo}
\end{equation}
Maximise the ELBO over both $\theta$ and $\phi$ jointly with SGD.
The first term is a \emph{reconstruction} loss; the second is a
\emph{regulariser} pulling the posterior toward the prior. The only
trick to make this trainable is the
\textbf{reparameterisation trick}: write
$z = \mu_\phi(x) + \sigma_\phi(x) \odot \epsilon$ with
$\epsilon \sim \mathcal{N}(0,I)$ so that the stochasticity is
isolated in $\epsilon$ and gradients flow through $\mu$ and
$\sigma$.

For two Gaussians $\mathcal{N}(\mu,\sigma^2 I)$ and
$\mathcal{N}(0,I)$ in $\R^d$,
\begin{equation}
  \KL\!\big(\mathcal{N}(\mu,\sigma^2 I) \,\|\, \mathcal{N}(0,I)\big) \;=\; \tfrac{1}{2}\sum_{i=1}^d \big( \sigma_i^2 + \mu_i^2 - 1 - \log\sigma_i^2 \big).
\end{equation}
Memorise this; you will reuse it in every VAE you ever write.

\subsubsection{Generative Adversarial Network (GAN)}

A \emph{generator} $G_\theta : \R^d \to \mathcal{X}$ maps noise
$z \sim p_z$ to samples $G(z)$; a
\emph{discriminator} $D_\psi : \mathcal{X} \to (0,1)$ tries to
tell real samples from generated ones. The \textbf{min-max} game is
\begin{equation}
  \min_\theta \max_\psi \; \E_{x \sim p_\text{data}}\log D_\psi(x) + \E_{z \sim p_z} \log\!\big(1 - D_\psi(G_\theta(z))\big).
\end{equation}
At the optimum (with infinite capacity), $D$ recovers the
density ratio and $G$ produces samples from $p_\text{data}$. In
practice, training is unstable; the standard fixes are
\textbf{WGAN-GP} (Wasserstein objective + gradient penalty),
\textbf{spectral normalisation} of $D$, \textbf{TTUR}
(separate learning rates), and the generator EMA. Conditional
variants (pix2pix, CycleGAN) condition on a second image and are the
GAN of choice for medical cross-modality synthesis.

\subsubsection{Diffusion models (in one paragraph)}

A diffusion model is a Markov chain that progressively adds
Gaussian noise to data over $T$ steps; a neural network learns the
\emph{reverse} step (the denoiser). Trained well, the reverse chain
samples from $p_\text{data}$. They have largely replaced GANs for
image synthesis. Their use as a \emph{prior} for medical-image
reconstruction (e.g.\ Chung \& Ye 2022) is currently one of the
hottest research areas.

\subsubsection{The thread tying it all together}

Look back at \eqref{eq:map-recon}. Anywhere $R(x)$ appears, you can
plug in:
\begin{itemize}[leftmargin=*,nosep]
  \item a hand-designed MRF energy (Week 2);
  \item the negative log-density of a VAE;
  \item the score (gradient of log-density) of a diffusion model.
\end{itemize}
The mathematical scaffolding of weeks 2 and 3 \emph{is} the same one
that frames the most modern reconstruction methods of 2024--26.

\subsection{Vocabulary you will meet}

\begin{description}[leftmargin=*]
  \item[\textbf{Latent variable.}] An unobserved random variable $z$
        that \emph{generates} the observation $x$.
  \item[\textbf{ELBO.}] Eq.~\eqref{eq:elbo}. The thing you maximise.
  \item[\textbf{Reparameterisation.}] Trick to backprop through
        $z \sim q_\phi(z|x)$ by writing $z = g_\phi(x, \epsilon)$.
  \item[\textbf{Posterior collapse.}] A failure mode of VAEs where
        the decoder ignores $z$ and the encoder outputs the prior.
  \item[\textbf{$\beta$-VAE.}] Multiply the KL term by $\beta$ to
        trade reconstruction for disentanglement.
  \item[\textbf{Mode collapse.}] GAN failure where $G$ produces
        only one (or a few) outputs.
  \item[\textbf{WGAN-GP.}] Wasserstein loss with a gradient penalty
        for stable adversarial training.
  \item[\textbf{Spectral normalisation.}] Constrain the spectral
        norm of $D$'s weights to enforce Lipschitzness.
  \item[\textbf{Pix2pix / CycleGAN.}] Paired and unpaired
        image-to-image translation GANs.
  \item[\textbf{DDPM.}] Denoising Diffusion Probabilistic Model
        (Ho, Jain \& Abbeel 2020).
  \item[\textbf{Score function.}] $\nabla_x \log p(x)$; what a
        diffusion / score-based model approximates.
\end{description}

\subsection{Pointers}

Lilian Weng's blog is the best single source of explainers for VAEs
and GANs. The original Kingma--Welling (VAE) and Goodfellow (GAN)
papers are short and worth reading in full. For diffusion, Lilian
Weng's blog again, plus the DDPM paper. For the medical application
of generative priors, Chung \& Ye (2022) is the cleanest read.

% =============================================================
\section{Weeks 6.5--8 (\texttt{week\_7}), part 1 --- Paper reading + guided implementation}
% =============================================================

\subsection{Big picture}

You now have the vocabulary of weeks 1--6. The point of week 7 is to
\emph{prove it} by reading a real research paper and reproducing
one experiment from it. The strongest signal that you understand a
method is not that you can summarise it --- it's that you can rebuild
it.

\subsubsection{Reading method}

Use Andrew Ng's three-pass method:
\begin{enumerate}[nosep, leftmargin=*]
  \item Pass 1 (5--10 min): title, abstract, figures, conclusions.
        Decide whether to keep going.
  \item Pass 2 ($\sim$1 h): read carefully but skip proofs. Write a
        one-page summary covering: problem, prior work and its
        limitation, key idea, method, experiments, results,
        limitations.
  \item Pass 3 (a day): re-implement the key equation or one
        experiment. Where you get stuck is exactly where you don't
        understand it yet.
\end{enumerate}

\subsubsection{Choosing a paper}

Pick something where:
\begin{itemize}[leftmargin=*,nosep]
  \item the problem connects to weeks 2--6 (don't go to a totally
        new sub-field this week);
  \item there is \emph{public code} and a \emph{public dataset};
  \item the paper is \emph{at most 12 pages};
  \item the venue is MICCAI, MedIA, TMI, IPMI, MIDL, NeurIPS.
\end{itemize}
A starter pool is in \texttt{week\_7/Readme.md}; Prof.\ Awate's
\texttt{Slides\_ProjectTopicExamples.pdf} is another good source.

\subsection{Deliverables}

A one-page summary, hand-written derivations of the central
equations, a reproduction notebook, and a short note describing the
gap between the paper's claim and your reproduction.

% =============================================================
\section{Weeks 6.5--8 (\texttt{week\_8}), part 2 --- Final project}
% =============================================================

\subsection{Deliverables}

A 6--10 page report, 10--15 talk slides, a clean reproducible
\texttt{code/} folder (with \texttt{requirements.txt}), a
\texttt{results/} folder, and a short live demo.

\subsection{Grading dimensions (mirroring CS-736)}

\begin{itemize}[leftmargin=*,nosep]
  \item Difficulty of the chosen topic.
  \item Depth of theoretical reading and understanding.
  \item Depth of empirical analysis: have you compared against a
        baseline? Have you run on a real test set with appropriate
        metrics?
  \item Quality of the slides.
  \item Reproducibility of the code.
\end{itemize}

\subsection{Suggested project shapes}

\begin{itemize}[leftmargin=*,nosep]
  \item Bayesian denoising with a learned prior --- CNN denoiser vs.
        MRF MAP baseline.
  \item Under-sampled MRI reconstruction --- model-based unrolled
        network (e.g.\ MoDL) on fastMRI vs.\ zero-filled baseline.
  \item Brain-tumour segmentation --- U-Net on BraTS vs.\
        nnU-Net out-of-the-box.
  \item Shape-prior segmentation --- ASM (week 4) combined with a
        CNN initialiser.
  \item Cross-modality synthesis --- CycleGAN MR$\leftrightarrow$CT
        on a small paired dataset.
  \item VAE anomaly detection on chest X-rays.
\end{itemize}

\subsection{Common mistakes to avoid}

No baseline. Training loss as the headline metric. Tiny test set
with no confidence interval. No qualitative failure analysis. Code
that won't run on someone else's machine.

% =============================================================
\section*{Closing note}
% =============================================================

\addcontentsline{toc}{section}{Closing note}

You have, over eight weeks, walked from ``what is a pixel'' to
``score-based diffusion priors for MRI''. The vocabulary you have now
is enough to read the proceedings of MICCAI cover-to-cover and to
join, contribute to, or lead, a medical-imaging research project.

Good luck. The work in this field is meaningful; many of the people
who do it make patients' lives demonstrably better. Take that
seriously.

\end{document}
